\section{설치 및 빌드 (Build \& Install)}

\subsection{현재 확인된 빌드 특성}
\texttt{include/oqs/oqsconfig.h} 기준:
\begin{itemize}[leftmargin=*]
\item 버전: \texttt{OQS\_VERSION\_TEXT = 0.15.0}
\item 플랫폼: \texttt{x86\_64 Linux}, 배포 빌드 플래그 활성
\item OpenSSL 연동: \texttt{OQS\_USE\_OPENSSL=1}
\item 스레드: \texttt{OQS\_USE\_PTHREADS=1}
\end{itemize}

\subsection{권장 빌드 절차}
프로젝트 루트에서 필요한 테스트 대상을 빌드한다.
예를 들어 KEM 테스트 바이너리는 다음 명령으로 생성한다.

\begin{lstlisting}[style=cstyle,caption={테스트 바이너리 빌드 예시}]
make kat_kem
\end{lstlisting}

SIG 테스트 바이너리가 필요하면 \texttt{make kat\_sig}를 실행한다.

\subsection{알고리즘 오브젝트 빌드 예시}
각 알고리즘 구현 파일은 다음과 같이 컴파일러를 직접 호출해 오브젝트 파일로 빌드할 수 있다.

\begin{lstlisting}[style=cstyle,caption={SLH-DSA 구현 파일 컴파일 예시}]
/usr/bin/cc -I./include \
    -std=gnu11 -fPIC -fvisibility=hidden -Wa,--noexecstack -O3 -fomit-frame-pointer -fdata-sections -ffunction-sections -Wl,--gc-sections -DMLD_CONFIG_PARAMETER_SET=44 -DMLD_CONFIG_FILE="../../integration/liboqs/config_c.h" \
    -MMD -MP -o build/sig/slh_dsa/wrappers/pure/slh_dsa_pure_shake_128f.o \
    -c ./src/sig/slh_dsa/wrappers/pure/slh_dsa_pure_shake_128f.c
\end{lstlisting}

\subsection{산출물 확인}
\begin{lstlisting}[style=cstyle,caption={주요 정적 라이브러리 산출물}]
build/lib/liboqs.a
build/lib/liboqs-internal.a
\end{lstlisting}

\subsection{테스트 실행 예시}
\begin{lstlisting}[style=cstyle,caption={KAT 기반 KEM 테스트 예시}]
./build/tests/kat_kem ml-kem-768 > /tmp/kat_mlkem768.out
echo $?
\end{lstlisting}
